\chapter{Hematology}

\medterm{ABO Incompatibility}-- A condition occurring when ABO blood group antigens on red blood cells are mismatched between a mother and fetus, leading to hemolytic disease of the fetus and newborn (HDFN). ABO incompatibility is the most common cause of HDFN, occurring more frequently than Rh incompatibility but generally causing milder disease. It most commonly affects neonates of blood group A or B born to group O mothers, because group O individuals produce IgG anti-A and anti-B antibodies that can cross the placenta. Unlike Rh disease, ABO incompatibility frequently affects first pregnancies and rarely causes severe hemolysis requiring exchange transfusion.

\medterm{Acute leukemia}-- leukemia characterized by a rapid onset, clonal proliferation of immature blast cells (>20\% of bone marrow cells), and rapid progression without treatment. Subtypes: acute lymphoblastic leukemia (ALL, B-cell or T-cell, most common childhood cancer, peak age 2-5) and acute myeloid leukemia (AML, more common in adults, median age 65-70). Clinical features: bone marrow failure (anemia--fatigue, pallor; thrombocytopenia--bruising, petechiae, bleeding; neutropenia--fever, recurrent infections), organ infiltration (lymphadenopathy, hepatosplenomegaly, bone pain, gum hypertrophy in AML M4/M5, skin infiltration). Emergency complications: hyperleukocytosis (blast count >100,000, risk of leukostasis, requiring urgent leukapheresis), tumor lysis syndrome, disseminated intravascular coagulation (especially in AML M3/promyelocytic). Diagnosis: full blood count and blood film, bone marrow aspirate and trephine biopsy, flow cytometry (immunophenotyping), cytogenetics (karyotype, FISH), and molecular genetics (NPM1, FLT3-ITD, CEBPA, RUNX1, BCR-ABL for ALL, PML-RARA for APL).

\medterm{Acute Lymphoblastic Leukemia}-- A clonal malignancy of lymphoid progenitor cells, most common in children (peak age 2-5 years) but also occurring in adults with a worse prognosis with increasing age. B-cell ALL (85 percent of childhood, 75 percent of adult ALL) is classified by immunophenotype: early pre-B (CD19+, CD10-, TdT+), common ALL (CD19+, CD10+, TdT+), pre-B (CD19+, cytoplasmic mu heavy chain+), and mature B-cell (Burkitt-type, CD19+, surface immunoglobulin+, CD10+). T-cell ALL (15\% of childhood, 25\% of adult ALL) presents with mediastinal mass and has a poorer prognosis.

\medterm{Acute Megakaryoblastic Leukemia}-- A rare subtype of acute myeloid leukemia characterized by the proliferation of megakaryoblasts in the bone marrow and peripheral blood. It accounts for approximately 1-2 percent of AML in adults but is the most common AML subtype in children with Down syndrome (trisomy 21), where it occurs in 1-5 percent of cases (called transient abnormal myelopoiesis in neonates, which often resolves spontaneously). AMKL is associated with extensive bone marrow fibrosis and a poor prognosis.

\medterm{Acute Promyelocytic Leukemia}-- A distinct subtype of acute myeloid leukemia (AML) characterized by the t(15;17)(q24;q21) translocation, which fuses the promyelocytic leukemia (PML) gene on chromosome 15 with the retinoic acid receptor alpha (RARA) gene on chromosome 17. This PML-RARA fusion protein blocks myeloid differentiation at the promyelocyte stage, leading to accumulation of abnormal promyelocytes in the bone marrow and peripheral blood. APL accounts for approximately 5-10 percent of AML cases.

\medterm{Anaemia}-- A condition in which the number of red blood cells or the hemoglobin concentration is below the normal range, reducing the oxygen-carrying capacity of the blood. Classified by etiology: blood loss (acute or chronic), decreased production (iron deficiency, vitamin B12/folate deficiency, bone marrow failure, anemia of chronic disease), and increased destruction (hemolytic anemias). Classified by morphology: microcytic (low MCV — iron deficiency, thalassemia, anemia of chronic disease), normocytic (normal MCV — acute blood loss, hemolysis, renal failure), macrocytic (high MCV — B12/folate deficiency, alcohol, MDS, hypothyroidism). Symptoms: fatigue, pallor, dyspnea, dizziness, tachycardia. Severe anemia can cause heart failure. Diagnosis: CBC, reticulocyte count, iron studies, vitamin B12/folate, and peripheral smear. Treatment: address underlying cause, iron supplementation, B12/folate replacement, erythropoietin, and blood transfusion for severe symptomatic anemia.

\medterm{Anemia}-- A reduction in the oxygen-carrying capacity of the blood, defined by a hemoglobin concentration below the normal range for age and sex: less than 13.5 g/dL in adult males, less than 12.0 g/dL in adult females. Anemia is classified by mechanism into three categories: decreased production (iron deficiency, vitamin B12 deficiency, folate deficiency, chronic disease, aplastic anemia, myelodysplastic syndrome), increased destruction (hemolytic anemias including hereditary spherocytosis, sickle cell disease, paroxysmal nocturnal haemoglobinuria), or blood loss.

\medterm{Antiphospholipid Syndrome}-- A systemic autoimmune thrombophilia characterized by arterial and venous thrombosis, pregnancy morbidity, and persistent antiphospholipid antibody positivity. Antiphospholipid antibodies are a heterogeneous group of autoantibodies targeting phospholipid-binding proteins: lupus anticoagulant (LA, detected by prolongation of phospholipid-dependent coagulation assays that fails to correct with mixing study, paradoxically associated with thrombosis), anticardiolipin antibodies (aCL, IgG or IgM, detected by ELISA), and anti-beta-2 glycoprotein I antibodies.

\medterm{Antithrombin Deficiency}-- A rare (1 in 2,000-5,000), strongly thrombophilic inherited disorder (autosomal dominant). Antithrombin (previously antithrombin III) is a serine protease inhibitor that inactivates thrombin and factors Xa, IXa, XIa, and XIIa. Heterozygous: 10-50 fold increased VTE risk, more severe than factor V Leiden. Presents with DVT and PE at a young age (<30), recurrent thrombosis, and unusual sites (mesenteric, cerebral veins). Diagnosis: low antithrombin activity (type I: reduced antigen and activity; type II: reduced activity). Heparin resistance is a clue. Management: anticoagulation for acute VTE; may require high-dose heparin; lifelong anticoagulation. Prophylaxis in high-risk situations is essential.

\medterm{Aplastic Anemia}-- A bone marrow failure syndrome characterized by pancytopenia (anemia, leukopenia, thrombocytopenia) and a hypocellular bone marrow with replacement by fat cells, resulting from immune-mediated destruction of hematopoietic stem cells. Acquired aplastic anemia accounts for 80-90 percent of cases, with autoimmune T-cell-mediated destruction of CD34+ hematopoietic progenitors as the primary mechanism, supported by response to immunosuppressive therapy. Causes include idiopathic (60-70\%), drugs, toxins, viruses, and pregnancy.

\medterm{Bleeding}-- The escape of blood from the vascular system, which may be external (visible hemorrhage from a wound, orifice, or natural opening) or internal (bleeding into a body cavity, organ, or tissue). Bleeding is classified by source: arterial bleeding (bright red, pulsatile, and difficult to control due to high pressure), venous bleeding (dark red, steady flow), and capillary bleeding (slow oozing from superficial wounds). The clinical significance of bleeding depends on the rate and volume of blood loss, as well as the location of bleeding, with hemorrhage in critical areas such as the intracranial cavity, pericardium, or airway being immediately life-threatening even with relatively small volumes of blood loss.

\medterm{Bleeding Disorders}-- A heterogeneous group of conditions characterized by abnormal hemostasis, resulting from defects in blood vessels, platelets, or coagulation factors. They are classified into three categories: vascular disorders (Ehlers-Danlos syndrome, Henoch-Schonlein purpura, hereditary hemorrhagic telangiectasia), platelet disorders (thrombocytopenia from immune destruction [immune thrombocytopenia], decreased production [aplastic anemia, chemotherapy], or sequestration [hypersplenism]; and platelet function disorders (von Willebrand disease, uraemia, aspirin, NSAIDs)), and coagulation disorders (hemophilia A and B, von Willebrand disease, factor deficiencies, vitamin K deficiency, liver disease, anticoagulant therapy, and DIC).

\medterm{Blood Group}-- The classification of blood based on the presence or absence of specific antigens (agglutinogens) on the surface of red blood cells. The ABO system has four groups: A (A antigen), B (B antigen), AB (both A and B), and O (neither A nor B). The Rh system: RhD-positive (85\% of population) or RhD-negative (15\%). Blood type compatibility is critical for transfusion: group O RhD-negative is the universal donor (no ABO/RhD antigens); group AB RhD-positive is the universal recipient (no ABO/RhD antibodies). In pregnancy, RhD-negative women carrying RhD-positive fetuses may develop anti-D antibodies, causing hemolytic disease of the newborn in subsequent pregnancies — prevented by antenatal and postnatal anti-D immunoglobulin (RhoGAM).

\medterm{Blood Transfusion}-- The administration of whole blood or blood components (packed red blood cells, platelets, fresh frozen plasma, cryoprecipitate) to restore or maintain adequate circulating blood volume and oxygen-carrying capacity. Indications: packed red blood cells for symptomatic anemia (hemoglobin less than 7-8 g/dL in stable patients; higher thresholds for patients with acute coronary syndrome, hemoglobin less than 10 g/dL, or ongoing hemorrhage); platelets for active bleeding with thrombocytopenia less than $50 \times 10^9$/L; FFP for coagulopathy with INR >1.5; cryoprecipitate for fibrinogen <100 mg/dL.

\medterm{Burkitt Lymphoma}-- A highly aggressive B-cell lymphoma characterized by MYC translocation (t(8;14) IGH-MYC, t(8;22) IGL-MYC, t(2;8) IGK-MYC) and a "starry sky" morphology (tingible body macrophages). Three variants: endemic (African, EBV+, jaw/facial bones), sporadic (abdominal, ileocecal), immunodeficiency-associated (HIV). Clinical: rapidly growing mass, LDH extremely high, TLS risk. Treatment: intensive short-course chemotherapy (CODOX-M/IVAC, hyper-CVAD, LMB) + rituximab + CNS prophylaxis with intrathecal methotrexate.

\medterm{Chronic leukemia}-- leukemia characterized by an insidious onset, clonal proliferation of relatively mature blood cells, and a more indolent clinical course. Subtypes: chronic lymphocytic leukemia (CLL, most common adult leukemia, median age 72, monoclonal B-cell proliferation CD5+, CD23+, often asymptomatic, managed with watchful waiting in early stage) and chronic myeloid leukemia (CML, characterized by Philadelphia chromosome t(9;22) BCR-ABL fusion, trilineage myeloproliferation, progresses through chronic phase to accelerated phase to blast crisis).

\medterm{Chronic Lymphocytic Leukemia}-- The most common leukemia in Western countries, characterized by the clonal proliferation of mature-appearing but functionally incompetent B lymphocytes that accumulate in blood, bone marrow, and lymphoid tissues. Diagnosis requires a monoclonal B-lymphocyte count of 5,000 per microliter or greater with a characteristic immunophenotype: CD5+, CD19+, CD23+, CD200+, with weak surface immunoglobulin expression and weak CD20. CLL has a highly variable clinical course, with some patients having indolent disease and others progressing rapidly.

\medterm{Chronic Myeloid Leukemia}-- A myeloproliferative neoplasm characterized by the Philadelphia chromosome [t(9;22)(q34;q11.2)], which creates the BCR-ABL1 fusion gene encoding a constitutively active tyrosine kinase. CML has three phases: chronic phase (most common at diagnosis, 85-90 percent, characterized by low blast count less than 10 percent, thrombocytosis, and basophilia), accelerated phase (increasing blasts 10-19 percent, basophilia greater than 20 percent, or additional cytogenetic abnormalities), and blast phase (>20\% blasts, acute leukemia-like, poor prognosis).

\medterm{Chronic Myelomonocytic Leukemia}-- A rare myelodysplastic/myeloproliferative neoplasm characterized by persistent monocytosis (absolute monocyte count greater than 1,000 per microliter and greater than 10 percent of white blood cells) in the peripheral blood and bone marrow, with dysplastic features and variable blast count. CMML bridges the categories of MDS and myeloproliferative neoplasm, with clinical features of both: cytopenias and dysplasia (MDS features) plus splenomegaly and monocytosis, and increased risk of transformation to AML (15-30\%).

\medterm{Coagulation Factor Deficiencies}-- Rare factor deficiencies: Factor VII (most common rare bleeding disorder, prolonged PT, normal aPTT, mucosal bleeding, hemarthrosis; recombinant FVIIa for treatment), Factor X (prolonged PT/aPTT), Factor XI (Ashkenazi Jews, prolonged aPTT, bleeding with surgery/trauma, FFP for replacement), Factor XIII (normal PT/aPTT, clot solubility test, umbilical stump bleeding, intracranial hemorrhage; FXIII concentrate), Fibrinogen (afibrinogenemia, hypofibrinogenemia, dysfibrinogenemia; cryoprecipitate, fibrinogen concentrate.

\medterm{Direct Oral Anticoagulants}-- Direct oral anticoagulants (DOACs) target specific coagulation factors: dabigatran (direct thrombin inhibitor, factor IIa), rivaroxaban, apixaban, edoxaban (direct factor Xa inhibitors). Indications: VTE treatment/prevention, stroke prevention in atrial fibrillation (non-valvular), post-orthopedic surgery prophylaxis. Advantages: fixed dosing, no routine monitoring, rapid onset/offset, fewer drug/food interactions (vs warfarin). Reversal: idarucizumab (anti-dabigatran monoclonal antibody), andexanet alfa (anti-Xa reversal). DOACs are preferred over warfarin for most indications.

\medterm{Disseminated Intravascular Coagulation}-- A pathological state of simultaneous systemic activation of coagulation and fibrinolysis, leading to widespread microvascular thrombosis that causes organ ischemia and dysfunction, paradoxically accompanied by consumptive coagulopathy and increased bleeding risk. DIC is triggered by conditions that expose tissue factor to the circulation: sepsis (most common trigger, especially gram-negative and fungal), trauma (tissue factor release), malignancy (particularly adenocarcinomas and acute promyelocytic leukemia),.

\medterm{Dyscrasia}-- An historic term broadly referring to any abnormal condition of the blood or bone marrow, now used primarily in the context of plasma cell dyscrasias (also called monoclonal gammopathies). These include multiple myeloma, monoclonal gammopathy of undetermined significance (MGUS), Waldenstrom macroglobulinemia, AL amyloidosis, and POEMS syndrome. The term indicates a disorder of blood cell production or function, often involving the overproduction of a monoclonal immunoglobulin (M protein) by a clonal plasma cell population. Diagnosis: serum protein electrophoresis, immunofixation, serum free light chain assay, and bone marrow biopsy.

\medterm{Erythropoietin}-- A glycoprotein hormone (EPO) that regulates erythropoiesis--the production of red blood cells in the bone marrow. Produced primarily by interstitial peritubular cells in the kidney (~90\%) and hepatocytes (~10\%) in response to tissue hypoxia. EPO binds to EPOR receptors on erythroid progenitor cells in the bone marrow, promoting their survival, proliferation, and differentiation into mature erythrocytes. Deficiency causes anemia of chronic kidney disease (the most clinically significant EPO deficiency). Recombinant human EPO (epoetin alfa, darbepoetin alfa) is used to treat anemia in CKD, chemotherapy-induced anemia, and before elective surgery to reduce transfusion needs. Adverse effects: hypertension, thrombosis (especially when hemoglobin targets are too high), and pure red cell aplasia (rare, antibody-mediated, associated with subcutaneous administration of certain epoetin formulations).

\medterm{Essential Thrombocythemia}-- A chronic myeloproliferative neoplasm characterized by sustained platelet count elevation (greater than 450,000 per microliter) due to clonal megakaryocytic proliferation, with a risk of both thrombotic and hemorrhagic complications. The JAK2 V617F mutation is present in 50-60 percent of cases, CALR mutations in 20-30 percent, and MPL mutations in 3-5 percent; approximately 10-15 percent are triple-negative. Clinical presentation ranges from asymptomatic (incidentally found on routine blood count) to thrombotic events (DVT, PE, microvascular occlusions including erythromelalgia and TIAs) and hemorrhagic complications (GI bleeding, epistaxis).

\medterm{Factor V Leiden}-- The most common inherited thrombophilia, caused by a point mutation in the factor V gene (G1691A, R506Q) resulting in resistance to activated protein C (APC resistance). Heterozygous: 5-fold increased VTE risk; homozygous: 50-80 fold increased risk. Most common in Caucasians (3-7\% prevalence). Clinical significance: increased risk of DVT, PE, and recurrent pregnancy loss. Less clear association with arterial thrombosis. Management: anticoagulation for VTE (consider DOACs); asymptomatic individuals generally do not require prophylaxis except in high-risk situations (surgery, pregnancy, prolonged immobilization).

\medterm{Folate Deficiency}-- Folate (vitamin B9) deficiency causes megaloblastic anemia identical to B12 deficiency but without neurologic manifestations (except possible cognitive changes). Dietary sources: leafy greens, legumes, liver; daily requirement 400 mcg. Absorption in jejunum. Causes: inadequate intake (alcoholism, malnutrition, elderly), malabsorption (celiac, short bowel), increased demand (pregnancy, hemolysis, dialysis), drugs (methotrexate, trimethoprim, phenytoin, sulfasalazine). Lab: macrocytic anemia, hypersegmented neutrophils on peripheral smear, low folate levels.

\medterm{Folate-deficiency anemia}-- A type of megaloblastic anemia caused by insufficient folic acid, a B vitamin essential for DNA synthesis, cell division, and red blood cell production. Folate deficiency impairs thymidylate synthesis, causing defective DNA replication in rapidly dividing cells of the bone marrow, resulting in megaloblastic erythropoiesis with large, abnormally shaped erythrocyte precursors. Causes include inadequate dietary intake (most common in alcohol use disorder, older adults, fad diets, and anorexia), increased requirements (pregnancy, hemolysis, malignancy, chronic inflammatory disorders), and impaired absorption (short bowel, celiac disease, gastric bypass). Diagnosis confirmed by low serum folate and elevated homocysteine; methylmalonic acid is normal (distinguishes from B12 deficiency).

\medterm{Follicular Lymphoma}-- The most common indolent non-Hodgkin lymphoma, accounting for 20-25 percent of NHL cases, characterized by t(14;18) translocation leading to BCL2 overexpression and follicular growth pattern. Clinical presentation: painless lymphadenopathy, often widespread, with bone marrow involvement in 60-70 percent. Grading: 1-2 (follicular, indolent), 3A (centrocytes + centroblasts), 3B (centroblasts only, aggressive). FLIPI (Follicular Lymphoma International Prognostic Index) incorporates five adverse prognostic factors: age >60, serum LDH > normal, number of nodal sites >4, hemoglobin <12 g/dL, and Ann Arbor stage III-IV.

\medterm{G6PD Deficiency}-- Glucose-6-phosphate dehydrogenase (G6PD) deficiency is an X-linked enzymopathy affecting the pentose phosphate pathway, leading to impaired NADPH production and reduced glutathione regeneration, rendering erythrocytes susceptible to oxidative hemolysis. Common variants: African A- (mild), Mediterranean (severe), Mahidol (SE Asia). Triggers: fava beans, oxidant drugs (primaquine, dapsone, sulfonamides, rasburicase, methylene blue), infections. Clinical: acute hemolytic anemia (back pain, dark urine (haemoglobinuria), jaundice, anemia); self-limiting within 1-2 weeks of trigger removal.

\medterm{Granulocyte}-- A type of white blood cell (leukocyte) containing cytoplasmic granules that are visible on light microscopy. Granulocytes are produced in the bone marrow (myelopoiesis, 3-4 days maturation) and are part of the innate immune system. Three types: (1) Neutrophils (40-70\% of WBCs): bilobed nucleus, pale pink granules, primary defence against bacteria and fungi, short half-life (6-8 hours). (2) Eosinophils (1-4\%): bilobed nucleus, orange-red granules, defence against parasites and involvement in allergic inflammation. (3) Basophils (<1\%): bilobed nucleus, large dark purple granules, release histamine and heparin in allergic reactions. Granulocyte count (total): 2.0--7.5 $\times$ 10$^9$/L. Neutropenia (ANC <1.5) increases infection risk; agranulocytosis (ANC <0.5) is a medical emergency.

\medterm{Haemangioma}-- A benign vascular tumor composed of proliferating endothelial cells, the most common soft tissue tumor of infancy. Haemangiomas present shortly after birth (not usually visible at birth), undergo a proliferative phase (first 6-12 months, rapid growth), followed by involution (slow regression over 5-10 years with replacement by fibrofatty tissue). 60\% occur in the head and neck region. Most are superficial (strawberry naevus), but can be deep (blue-hued) or mixed. Complications: ulceration (10-20\%, most common), bleeding, infection, visual axis obstruction (periocular), airway obstruction (subglottic), Kasabach-Merritt phenomenon (platelet trapping, consumptive coagulopathy). Treatment options: observation (majority involute spontaneously), topical timolol for superficial lesions, oral propranolol (2-3 mg/kg/day, first-line for problematic lesions), laser therapy for residual telangiectasia after involution, surgical excision for complications or residual deformity.

\medterm{Haemarthrosis}-- Bleeding into a joint cavity. Causes: trauma (most common), hemophilia (spontaneous haemarthrosis is characteristic, especially in severe deficiency of factor VIII or IX, common in knees, elbows, ankles), anticoagulant therapy, pigmented villonodular synovitis (PVNS, recurrent haemarthrosis), joint surgery, and vascular malformations. In hemophilia, recurrent haemarthrosis leads to synovitis, cartilage damage, joint space narrowing, and ultimately haemophilic arthropathy (chronic pain, stiffness, deformity, ankylosis). Management: factor replacement (aim for 50-100\% activity), joint aspiration only if tense and painful (risk of infection), rest, ice, splinting, physiotherapy to prevent contractures.

\medterm{Haematocrit}-- The proportion of blood volume occupied by red blood cells (RBCs), expressed as a percentage or decimal fraction. Normal reference range: men 42-50\%, women 37-47\% (varies by altitude, age, and laboratory). Haematocrit is measured by centrifugation of whole blood in a microhaematocrit tube (PCV--packed cell volume) or calculated from mean corpuscular volume (MCV) and RBC count (Hct = MCV x RBC count / 10). Clinical interpretation: low haematocrit (anemia--decreased RBC production, increased destruction, or blood loss), high haematocrit (polycythemia--primary: polycythemia vera; secondary: chronic hypoxia, high altitude, COPD, cyanotic heart disease, EPO-secreting tumors, smoking). The haematocrit is used to assess the severity of anemia and to guide transfusion decisions (trigger typically <21-24\% or <7-8 g/dL hemoglobin depending on clinical context).

\medterm{Haematologist}-- A medical doctor who specializes in the diagnosis, treatment, and prevention of diseases of the blood and blood-forming organs. Conditions managed: anemia, bleeding disorders (hemophilia, von Willebrand disease), thrombophilia, hematologic malignancies (leukemia, lymphoma, myeloma, MDS, MPN), and hemoglobinopathies (sickle cell disease, thalassemia). Hematologists interpret laboratory tests (CBC, coagulation studies, bone marrow biopsy, flow cytometry, cytogenetics, molecular testing) and administer treatments including chemotherapy, targeted therapy, immunotherapy, anticoagulation, clotting factor replacement, and stem cell transplantation.

\medterm{Haematoma}-- A localised collection of blood outside blood vessels, usually clotted, within an organ or tissue. Haematomas form when trauma or disease causes a blood vessel to rupture. They are classified by location: subcutaneous (ecchymosis/bruise), intramuscular, subdural, epidural, intracerebral, subungual (under nail), perianal, rectus sheath, psoas, septal (nasal), subchorionic/retroplacental (pregnancy), and retroperitoneal. A haematoma may expand (especially in anticoagulated patients) and can cause complications: compression of adjacent structures (nerve compression, compartment syndrome), infection if it becomes secondarily infected, organisation and fibrosis, or calcification. Management depends on size and location: most small haematomas resorb spontaneously over weeks; large or expanding haematomas may require drainage (surgical or percutaneous) and reversal of any coagulopathy.

\medterm{Haematuria}-- The presence of blood in the urine. Classification: gross (visible to the naked eye, indicates bleeding beyond the nephron) vs microscopic (detected only on urinalysis or microscopy, >3 RBCs per high-power field). Aetiologies: (1) Glomerular: IgA nephropathy (most common cause of macroscopic haematuria), post-streptococcal GN, membranoproliferative GN, Alport syndrome, thin basement membrane nephropathy. Glomerular haematuria often presents with dysmorphic RBCs (acanthocytes) and RBC casts. (2) Non-glomerular: urinary tract infection (most common cause), urolithiasis (renal calculi), malignancy (bladder cancer, RCC, prostate cancer), benign prostatic hyperplasia (BPH), trauma, exercise-induced haematuria, anticoagulation. (3) Factitious: menstrual contamination, factitious haematuria. Evaluation: history (timing, associated symptoms, smoking history, anticoagulant use), urinalysis and culture, urine cytology, imaging (CT urogram or ultrasound), cystoscopy (for persistent or unexplained haematuria, especially in older adults).

\medterm{Haemodialysis}-- A medical procedure that filters waste products, toxins, and excess fluid from the blood when the kidneys are unable to do so (end-stage renal disease, ESRD). Mechanism: blood is pumped through an extracorporeal circuit and passes through a dialyser (artificial kidney) containing thousands of hollow fibers with a semipermeable membrane. Dialysis accomplishes: diffusion (removal of small solutes--urea, creatinine, potassium, phosphate down a concentration gradient), ultrafiltration (removal of excess fluid by applying hydrostatic pressure), and convection (removal of middle molecules by solvent drag). Access: arteriovenous fistula (AVF, gold standard), arteriovenous graft (AVG), or central venous catheter (tunnelled, non-tunnelled). Typical schedule: 3-4 hours, 3 times per week. Complications: hypotension (most common, due to ultrafiltration), muscle cramps, disequilibrium syndrome, air embolism, access complications (stenosis, thrombosis, aneurysm, infection), dialyser reactions (hypersensitivity to ethylene oxide or heparin), amyloidosis (beta-2 microglobulin accumulation).

\medterm{Haemophilia}-- An X-linked recessive bleeding disorder caused by deficiency of coagulation factor VIII (hemophilia A, 80-85\%) or factor IX (hemophilia B, Christmas disease, 15-20\%). Incidence: 1 in 5,000 male births (hemophilia A), 1 in 30,000 (hemophilia B). Severity correlates with residual factor activity: severe (<1\% activity, spontaneous bleeding into joints and muscles), moderate (1-5\%, bleeding after minor trauma), mild (6-40\%, bleeding after major trauma or surgery). Clinical features: spontaneous haemarthrosis (knees, elbows, ankles--most characteristic), intramuscular haematomas (iliopsoas, calf, forearm), easy bruising, prolonged bleeding after dental extractions or surgery, haematuria, intracranial hemorrhage (most common cause of death). Diagnosis: prolonged aPTT with normal PT and platelet count, factor VIII and IX assays (low). Treatment: factor replacement (recombinant or plasma-derived, on-demand or prophylaxis), desmopressin (DDAVP, releases von Willebrand factor, effective in mild hemophilia A), emicizumab (bispecific monoclonal antibody, mimics factor VIII function, subcutaneous prophylaxis for hemophilia A). Complications: development of inhibitors (alloantibodies against factor VIII/IX, 20-30\% of severe hemophilia A, renders standard replacement ineffective), haemophilic arthropathy (chronic joint damage), transfusion-transmitted infections (HIV, hepatitis B/C--now rare due to viral inactivation and recombinant products).

\medterm{Haemoptysis}-- Coughing up blood or blood-stained sputum from the lower respiratory tract. Aetiologies: bronchitis (most common, 25-50\%, acute or chronic), bronchiectasis (dilated airways, chronic inflammation), lung cancer (squamous cell most commonly, 20\% present with haemoptysis), tuberculosis (especially cavitary or endobronchial), pulmonary embolism (haemoptysis with pleuritic chest pain and dyspnoea), pneumonia (especially Klebsiella--currant jelly sputum), lung abscess, bronchiectasis, aspergilloma (fungus ball in a cavity), mitral stenosis (pulmonary venous hypertension causing rupture of bronchial varices), iatrogenic (bronchoscopy, lung biopsy, anticoagulation), and vasculitides (Wegener granulomatosis, Goodpasture syndrome). Massive haemoptysis (>200 mL/24 hours or >50 mL/hour) is life-threatening due to asphyxiation rather than exsanguination. Diagnostic workup: chest X-ray, CT chest, bronchoscopy, sputum culture and cytology, complete blood count, coagulation profile. Management: ABC, position patient with bleeding side down, bronchial artery embolisation (BAE) for massive haemoptysis, surgery (lobectomy, pneumonectomy) for refractory cases.

\medterm{Hairy Cell Leukemia}-- A rare, indolent B-cell chronic lymphoproliferative disorder characterized by the accumulation of mature B lymphocytes with characteristic cytoplasmic projections ("hairy cells") in the bone marrow, peripheral blood, and spleen. It accounts for less than 2 percent of adult leukemias and predominantly affects middle-aged men (male-to-female ratio approximately 4-5:1). Classic HCL is driven by the BRAF V600E mutation, present in virtually all cases. Clinical presentation includes pancytopenia, splenomegaly (often massive), recurrent infections, and fatigue.

\medterm{Hematocrit}-- The percentage of total blood volume composed of packed red blood cells, one of the most commonly ordered laboratory tests in clinical medicine. The hematocrit is measured by centrifugation of anticoagulated whole blood in a capillary tube or automated analyzer, separating the cellular elements from plasma; the normal range is 41-53 percent in adult men and 36-46 percent in adult women, with lower values in children and during pregnancy (physiologic hemodilution). The hematocrit reflects both the packed cell volume and the plasma volume, so it is influenced by both the red blood cell mass and the plasma volume, with a low haematocrit indicating anemia and a high haematocrit suggesting polycythemia.

\medterm{hematology}-- The branch of medicine concerned with the study, diagnosis, treatment, and prevention of disorders of the blood and blood-forming organs (bone marrow, lymph nodes, spleen). Subspecialties: benign hematology (anemia, haemoglobinopathies--sickle cell disease, thalassemia, hemostatic disorders--hemophilia, von Willebrand disease, immune thrombocytopenia), malignant hematology (acute and chronic leukemias, myelodysplastic syndromes, myeloproliferative neoplasms, lymphoma, multiple myeloma, amyloidosis), transfusion medicine (blood banking, donor management, component therapy, apheresis, stem cell transplantation), haemato-oncology (chemotherapy, targeted therapy, immunotherapy, CAR-T cells, HSCT). Diagnostic tools: full blood count (FBC), blood film, bone marrow aspiration and trephine biopsy, flow cytometry, cytogenetics, molecular genetics (PCR, NGS), coagulation studies, hemoglobin electrophoresis, serum protein electrophoresis.

\medterm{Hemoglobin}-- The oxygen-carrying metalloprotein in red blood cells, composed of four globin subunits each containing an iron-porphyrin (heme) group. Adult hemoglobin (HbA): alpha2-beta2 tetramer (two alpha chains, two beta chains). Fetal hemoglobin (HbF): alpha2-gamma2, higher oxygen affinity for oxygen transfer across the placenta. hemoglobin reversibly binds oxygen (4 O2 molecules per tetramer) through cooperative binding (sigmoidal oxygen dissociation curve, p50 = 26 mmHg). The Bohr effect: decreased pH and increased CO2 shift the curve rightward, increasing oxygen unloading in tissues. Normal: men 13.5-17.5 g/dL, women 12.0-16.0 g/dL. Abnormal haemoglobins: sickle cell disease (HbS, Glu6Val in beta-globin), HbC, HbE, metHb, carboxyHb (CO poisoning, cherry-red appearance). HbA1c: glycated hemoglobin, measures average blood glucose over 2-3 months, used for diabetes monitoring. Buffy coat: anticoagulated blood centrifuged shows three layers: plasma (top), buffy coat (middle, WBCs and platelets), packed RBCs (bottom).

\medterm{Hemoglobinopathies: Thalassemias}-- Thalassemias are inherited disorders of globin chain synthesis imbalance: alpha-thalassemia (HBA1/HBA2 deletions) and beta-thalassemia (HBB mutations). Beta-thalassemia major (transfusion-dependent): severe anemia presenting in infancy, ineffective erythropoiesis, marrow expansion, extramedullary hematopoiesis, iron overload from transfusions. Management: chronic transfusions (maintain Hb >9-10 g/dL), iron chelation (deferasirox, deferoxamine, deferiprone), splenectomy if hypersplenism, luspatercept (activin receptor ligand trap for beta-thalassemia), and allogeneic HSCT as definitive therapy.

\medterm{hemolysis}-- The premature destruction of red blood cells, releasing hemoglobin into the plasma. Extravascular hemolysis (most common): RBCs are destroyed by macrophages in the spleen and liver — causes unconjugated hyperbilirubinemia, increased LDH, and decreased haptoglobin. Intravascular hemolysis: RBCs rupture within the circulation — causes hemoglobinemia, hemoglobinuria, and decreased haptoglobin. Causes: inherited (hereditary spherocytosis, G6PD deficiency, sickle cell disease, thalassemia) and acquired (autoimmune hemolytic anemia, microangiopathic hemolytic anemia — DIC, TTP, HUS; mechanical heart valves, infection, drugs, transfusion reaction). Lab findings: reticulocytosis, elevated LDH, unconjugated bilirubin, decreased haptoglobin, and abnormal peripheral smear (spherocytes, schistocytes, sickle cells). Treatment: address underlying cause, folic acid supplementation, corticosteroids for AIHA, and splenectomy for selected cases.

\medterm{Hemolytic Anemia}-- Anemia due to accelerated destruction of red blood cells (RBCs), classified as intrinsic (RBC defect) or extrinsic (external factors). Intrinsic: hereditary spherocytosis, G6PD deficiency, sickle cell disease, thalassemia. Extrinsic: autoimmune hemolytic anemia (warm or cold agglutinin), microangiopathic hemolytic anemia (TMA, DIC, TTP/HUS), drug-induced, and mechanical (prosthetic valves). Laboratory: elevated LDH, indirect bilirubin, reticulocyte count; decreased haptoglobin; positive direct Coombs test in immune-mediated. Treatment depends on cause: corticosteroids for warm AIHA, splenectomy, and immunosuppressants.

\medterm{Hemolytic Anemias: Autoimmune}-- Autoimmune hemolytic anemia (AIHA): warm (IgG, 70-80\%, extravascular hemolysis in spleen, DAT positive for IgG +/- C3, associated with CLL, SLE, drugs, idiopathic), cold agglutinin disease (IgM, complement-mediated, intravascular/extravascular, DAT positive for C3 only, associated with lymphoproliferative, Mycoplasma), mixed (both warm and cold), paroxysmal cold hemoglobinuria (Donath-Landsteiner antibody, biphasic hemolysin, post-viral in children). Drug-induced (penicillins, cephalosporins, methyldopa). Treatment: corticosteroids (prednisone 1 mg/kg/day) for warm AIHA; rituximab, splenectomy, IVIG, and immunosuppressants for refractory cases; avoid transfusions if possible; keep patient warm for cold agglutinin disease.

\medterm{Hemolytic Anemias: Hereditary}-- Hereditary spherocytosis (HS): cytoskeletal protein defects (ankyrin, band 3, protein 4.2, alpha/beta spectrin) -> spherical RBCs, splenic trapping, hemolysis. AD (75\%) or AR. Clinical: anemia, jaundice, splenomegaly, gallstones. Spherocytes on smear, increased MCHC, osmotic fragility test, eosin-5-maleimide (EMA) binding (flow cytometry). Treatment: folic acid, splenectomy (if moderate/severe, after age 5-6, pneumococcal/meningococcal/Hib vaccines pre-op), partial splenectomy. Hereditary elliptocytosis: milder, usually asymptomatic; hereditary pyropoikilocytosis: severe hemolysis.

\medterm{hemolytic Disease of the Newborn}-- A condition in which maternal antibodies destroy fetal/neonatal red blood cells (RBCs), causing hemolytic anemia, hyperbilirubinaemia, and in severe cases hydrops fetalis. Most commonly due to Rh incompatibility: an RhD-negative mother carries an RhD-positive fetus; fetomaternal hemorrhage at delivery (or during pregnancy) sensitises the mother to produce anti-D IgG antibodies; subsequent pregnancies with RhD-positive fetuses are vulnerable. ABO incompatibility (mother O, baby A or B) is more common but typically milder. Pathophysiology: maternal IgG crosses the placenta, coats fetal RBCs, and causes extravascular hemolysis in the fetal spleen. Severity ranges from mild jaundice to severe anemia, hepatosplenomegaly, hydrops (anasarca, ascites, pleural effusion, heart failure), and intrauterine death (hydropic). Diagnosis: maternal antibody screening (indirect Coombs test), paternal zygosity testing, fetal genotyping (cell-free fetal DNA), serial MCA-PSV Doppler ultrasound to detect fetal anemia. Management: intrauterine transfusion (IUT) for severe anemia, early delivery, phototherapy for neonatal jaundice, exchange transfusion for severe hyperbilirubinaemia (neurotoxicity--kernicterus). Prevention: routine antenatal anti-D prophylaxis (RAADP) at 28 weeks and postpartum anti-D IgG within 72 hours of delivery for all non-sensitised RhD-negative women. Screening and prophylaxis have reduced HDN incidence by 90\%.

\medterm{Hemolytic Jaundice}-- Jaundice resulting from excessive hemolysis (breakdown of red blood cells) causing overproduction of unconjugated bilirubin that exceeds the liver's capacity for conjugation and excretion. See Jaundice for complete overview. Causes: hereditary spherocytosis, G6PD deficiency, sickle cell disease, thalassemia, autoimmune hemolytic anemia, hemolytic transfusion reaction, malaria, and microangiopathic hemolytic anemia (MAHA). Laboratory features: elevated unconjugated (indirect) bilirubin, increased reticulocyte count, low haptoglobin, elevated LDH, raised urobilinogen in urine, and evidence of hemolysis on blood film (spherocytes, schistocytes, bite cells).

\medterm{Hemolytic Uremic Syndrome}-- A thrombotic microangiopathy characterized by microangiopathic hemolytic anemia, thrombocytopenia, and acute kidney injury. Typical (Shiga toxin-producing E. coli, STEC-HUS): O157:H7 (most common), non-O157; follows diarrheal prodrome; children; supportive care (no antibiotics, may increase toxin release); eculizumab not routinely recommended. Atypical (aHUS): complement dysregulation (mutations in CFH, CFI, MCP/CD46, C3, CFB, THBD, or autoantibodies to CFH); requires eculizumab (anti-C5 monoclonal antibody) and plasma exchange. Management of both types includes supportive care with transfusion and plasma exchange.

\medterm{hemorrhoids (Piles)}-- Vascular cushions in the anal canal composed of arteriovenous anastomoses (normal structures present from birth), which become symptomatic when they bleed, prolapse, or thrombose. Anatomy: internal hemorrhoids (above the dentate line, 3 primary positions--left lateral, right anterior, right posterior, covered by columnar epithelium, visceral innervation--painless), external hemorrhoids (below the dentate line, covered by squamous epithelium, somatic innervation--painful if thrombosed). Grading of internal hemorrhoids: grade 1 (confined to anal canal, no prolapse), grade 2 (prolapses on defecation but reduces spontaneously), grade 3 (prolapses and requires manual reduction), grade 4 (irreducible, permanently prolapsed). Risk factors: chronic constipation/straining, prolonged sitting, pregnancy, obesity, low-fibre diet, cirrhosis (portal hypertension). Symptoms: painless rectal bleeding (bright red blood on toilet paper or dripping into the bowl, typically coating stools), prolapse, pruritus ani, mucoid discharge, thrombosis (acute severe pain, perianal lump). Diagnosis: clinical (inspection of external hemorrhoids, digital rectal examination, anoscopy, proctoscopy). Management: lifestyle (high-fibre diet, adequate hydration increased fibre, avoid straining), topical treatments (anaesthetics, steroids, astringents), rubber band ligation (grade 1-2), sclerotherapy, infrared coagulation, haemorrhoidectomy (Milligan-Morgan open, Ferguson closed, stapled, THD--transanal haemorrhoidal dearterialisation).

\medterm{hemostasis}-- The physiological process that stops bleeding at the site of vascular injury. Three phases: (1) Primary hemostasis: vasoconstriction (reflex, endothelin-mediated), platelet adhesion (to exposed subendothelial von Willebrand factor via platelet glycoprotein Ib/IX/V, to collagen via GPVI), platelet activation (shape change, granule release--ADP, thromboxane A2, serotonin), and platelet aggregation (fibrinogen bridges between platelet GPIIb/IIIa receptors), forming a platelet plug. (2) Secondary hemostasis: coagulation cascade, culminating in cross-linked fibrin stabilising the platelet plug. (3) Fibrinolysis: tissue plasminogen activator (tPA) converts plasminogen to plasmin, which degrades fibrin into D-dimers and FDPs. The coagulation cascade is classically described as intrinsic (surface contact--factors XII, XI, IX, VIII) and extrinsic (tissue factor--factor VII) pathways converging on the common pathway (factor X, V, II, I). In vivo, the cell-based model of hemostasis more accurately describes initiation on tissue factor-bearing cells, amplification on platelets, and propagation on platelet surfaces. Disorders: hemophilia (deficiency of VIII or IX), von Willebrand disease (deficiency or dysfunction of vWF), thrombocytopenia, platelet function disorders, anticoagulation therapy.

\medterm{Heparin-Induced Thrombocytopenia}-- An immune-mediated adverse drug reaction caused by IgG antibodies against platelet factor 4 (PF4)-heparin complexes, leading to platelet activation, thrombocytopenia, and a prothrombotic state. HIT typically occurs 5-14 days after heparin exposure (unfractionated heparin > low molecular weight heparin). Clinical: thrombocytopenia (fall >50\% from baseline, nadir rarely <20k), thrombosis (venous: DVT/PE; arterial: limb ischemia, stroke, MI) in 20-50\%, skin necrosis at injection sites (50\%), and adrenal hemorrhage. Diagnosis: 4Ts score, HIT ELISA (PF4-heparin ELISA), serotonin release assay (SRA). Treatment: stop heparin, start non-heparin anticoagulant (argatroban, bivalirudin, fondaparinux, danaparoid). Avoid warfarin until platelets >150 (risk of venous limb gangrene).

\medterm{Hereditary Spherocytosis}-- The most common inherited hemolytic anemia in Northern Europeans (autosomal dominant 75\%, recessive 25\%). Caused by defects in RBC membrane proteins (spectrin, ankyrin, band 3, protein 4.2), leading to loss of membrane surface area, reduced deformability, and premature destruction in the spleen. Presents with hemolytic anemia, jaundice, splenomegaly, and reticulocytosis. Lab findings: spherocytes on peripheral smear, increased MCHC, negative direct Coombs test, elevated osmotic fragility. Complications: aplastic crisis (parvovirus B19), hemolytic crisis, gallstones (pigment stones), and leg ulcers. Treatment: folic acid supplementation; splenectomy for severe anemia.

\medterm{Hodgkin Lymphoma}-- A B-cell lymphoma characterized by the presence of Reed-Sternberg cells (large, binucleate or multinucleate cells with owl-eye nucleoli) in a background of reactive inflammatory cells. Classical HL (cHL) accounts for 95 percent of cases and is subdivided into nodular sclerosis (most common, mediastinal predominance, young adults), mixed cellularity, lymphocyte-rich, and lymphocyte-depleted subtypes. Nodular lymphocyte-predominant HL (NLPHL, 5 percent) features lymphocyte-predominant (LP) cells (CD20+, CD15-, CD30-) instead of Reed-Sternberg cells. Staging with PET-CT. Treatment: ABVD (doxorubicin, bleomycin, vinblastine, dacarbazine) or escalated BEACOPP for advanced disease; involved-field radiation for bulky disease; brentuximab vedotin for relapsed/refractory CD30+ disease.

\medterm{Hodgkin's Lymphoma}-- See Hodgkin Lymphoma.

\medterm{Immune Thrombocytopenia}-- An autoimmune disorder characterized by isolated thrombocytopenia (platelet count <100 x10\textsuperscript{9}/L) due to autoantibody-mediated platelet destruction and impaired platelet production. Primary ITP (no identifiable cause) vs secondary (SLE, CLL, HIV, HCV, drugs). Clinical: petechiae, purpura, epistaxis, menorrhagia; intracranial hemorrhage rare (<1\%). Diagnosis: exclusion of other causes (CBC, blood smear, HIV/HCV, TSH, ANA, bone marrow if >60 or atypical features. First-line: corticosteroids (prednisone). Second-line: TPO receptor agonists (eltrombopag, romiplostim), rituximab, mycophenolate mofetil. Splenectomy reserved for refractory cases. Platelet transfusion for life-threatening bleeding.

\medterm{Iron Deficiency Anemia}-- The most common anemia worldwide, affecting over 1 billion people, resulting from insufficient iron for hemoglobin synthesis. Iron balance requires adequate dietary intake (heme iron from animal sources better absorbed than non-heme iron from plant sources), efficient absorption (duodenal and proximal jejunal enterocytes, regulated by hepcidin), and minimal losses (menstruation, GI bleeding, pregnancy). Causes include chronic blood loss (menorrhagia, GI bleeding from GI lesions), malabsorption, and increased requirements (pregnancy, growth). Lab: microcytic hypochromic anemia, low ferritin, low serum iron, low transferrin saturation, increased TIBC, increased RDW. Diagnosis: iron studies with ferritin <30 ng/mL. Treatment: oral iron (ferrous sulfate 325 mg daily); IV iron for intolerance, malabsorption, severe deficiency, or chronic blood loss. Response measured by reticulocyte count increase within 7-10 days and Hb rise of 1-2 g/dL over 3-4 weeks.

\medterm{Leucopenia}-- A decreased total white blood cell count below the normal reference range (<4.0 $\times$ 10$^9$/L). Leucopenia is often due to a decrease in neutrophils (neutropenia) or lymphocytes (lymphopenia). Causes: infection (viral--HIV, EBV, influenza, hepatitis; overwhelming bacterial sepsis--endotoxin-induced bone marrow suppression; typhoid), bone marrow failure (aplastic anemia, myelodysplasia, leukemia, myelofibrosis), chemotherapy/radiotherapy (most common cause of moderate-severe neutropenia), drugs (carbimazole, methimazole, clozapine, sulfonamides, immunosuppressants, antipsychotics, anticonvulsants), autoimmune (SLE, Felty syndrome--neutropenia with RA and splenomegaly), hypersplenism (cirrhosis, portal hypertension), nutritional deficiency (vitamin B12, folate, copper), and ethnic/benign variants (pseudoneutropenia, Afro-Caribbean and African descent, stable neutrophil count 1.0--2.5 $\times$ 10$^9$/L with no increased infection risk). Clinical significance: increased risk of infection, especially when absolute neutrophil count (ANC) falls <1.0 $\times$ 10$^9$/L. Severe neutropenia (ANC <0.5) requires urgent evaluation. Management: treat underlying cause, stop offending drug, G-CSF (filgrastim) for chemotherapy-induced neutropenia, antibiotic prophylaxis for severe chronic neutropenia.

\medterm{Leukemia}-- A group of hematological malignancies characterized by the clonal proliferation of abnormal white blood cells (blasts) in the bone marrow and peripheral blood. leukemia is classified by: (1) clinical course: acute (rapid onset, >20\% blasts in bone marrow) vs chronic (insidious onset, more differentiated cells). (2) cell lineage: lymphoblastic (T-cell or B-cell precursor) vs myeloid (granulocytic, monocytic, erythroid, megakaryocytic). The four main subtypes are: acute lymphoblastic leukemia (ALL, most common in children), acute myeloid leukemia (AML, more common in adults), chronic lymphocytic leukemia (CLL, most common adult leukemia in Western countries), and chronic myeloid leukemia (CML, characterized by Philadelphia chromosome t(9;22) BCR-ABL fusion).

\medterm{Leukocytosis}-- An increased total white blood cell count above the normal reference range (>11 $\times$ 10$^9$/L in adults). Causes: physiological (exercise, stress, pregnancy, labor), infectious (most common, especially bacterial--neutrophilia; viral--lymphocytosis; parasitic--eosinophilia), inflammatory (autoimmune disease, vasculitis, gout, pancreatitis), tissue necrosis (myocardial infarction, burns, trauma, surgery), malignancy (myeloproliferative neoplasms--CML, polycythemia vera; solid tumors with bone marrow involvement), medications (corticosteroids, lithium, G-CSF), hematologic (hemolytic anemia, hemorrhage), and post-splenectomy. leukocytosis is further characterized by the specific cell line increased: neutrophilia (neutrophils >7.5 $\times$ 10$^9$/L, bacterial infection, stress, steroids, CML), lymphocytosis (lymphocytes >4.0 $\times$ 10$^9$/L, viral infection, CLL, pertussis), monocytosis (monocytes >0.8 $\times$ 10$^9$/L, TB, endocarditis, CMML), eosinophilia (eosinophils >0.5 $\times$ 10$^9$/L, parasitic infection, asthma, allergy, eosinophilic granulomatosis with polyangiitis, drug reaction), and basophilia (basophils >0.1 $\times$ 10$^9$/L, CML, polycythemia vera). A normal peripheral blood smear and lack of abnormal cells distinguishes reactive leukocytosis from leukemia.

\medterm{Lymphoma (Hodgkin's and Non-Hodgkin's)}-- See Hodgkin Lymphoma and Non-Hodgkin Lymphoma.

\medterm{Mantle Cell Lymphoma}-- An aggressive B-cell lymphoma characterized by t(11;14) translocation leading to cyclin D1 overexpression and CCND1-IGH fusion. Clinical: older males, widespread lymphadenopathy, splenomegaly, GI involvement (lymphomatous polyposis), bone marrow, peripheral blood (leukemic phase). Ki-67 index and MIPI (MCL International Prognostic Index) stratify risk. Treatment: young/fit: intensive chemoimmunotherapy (R-DHAP, R-HyperCVAD, Nordic MCL2/3) + autologous stem cell transplant. Older/unfit: less intensive regimens (R-CHOP, R-bendamustine, VR-CAP) plus rituximab maintenance. No cure; median survival 5-7 years.

\medterm{Multiple Myeloma}-- A plasma cell neoplasm characterized by clonal proliferation of malignant plasma cells in the bone marrow, producing monoclonal immunoglobulin (M-protein) and causing end-organ damage defined by the CRAB criteria: hypercalcemia (calcium greater than 11 mg/dL), renal insufficiency (creatinine greater than 2 mg/dL), anemia (hemoglobin less than 10 g/dL), and bone lesions (lytic lesions on imaging). The International Myeloma Working Group (IMWG) diagnostic criteria also include clonal bone marrow plasma cells \(\ge\)10\% or biopsy-proven bony or extramedullary plasmacytoma, and any one or more of the CRAB features or one or more biomarker of malignancy (clonal plasma cell proliferation \(\ge\)60\%, involved:uninvolved serum free light chain ratio \(\ge\)100, or >1 focal lesion on MRI). Treatment: triplet regimens (VRd: bortezomib, lenalidomide, dexamethasone; or Dara-Rd for transplant-ineligible), followed by autologous stem cell transplant in eligible patients.

\medterm{Myelodysplastic Syndromes}-- A heterogeneous group of clonal hematopoietic stem cell disorders characterized by ineffective hematopoiesis, dysplastic morphology in one or more myeloid lineages, cytopenias, and risk of transformation to acute myeloid leukemia. MDS predominantly affects older adults (median age 70-75 years) and may be primary (de novo) or secondary (therapy-related from prior chemotherapy or radiation). Clinical presentation includes fatigue (anemia), infections (neutropenia), and bleeding (thrombocytopenia). IPSS-R (Revised International Prognostic Scoring System) stratifies risk based on cytopenias, blast percentage, and cytogenetics. Management: supportive care (transfusions, growth factors), lenalidomide for del(5q), luspatercept for ring sideroblasts, hypomethylating agents (azacitidine, decitabine), and allogeneic HSCT for high-risk disease.

\medterm{Myelofibrosis}-- A myeloproliferative neoplasm characterized by bone marrow fibrosis, extramedullary hematopoiesis, and JAK2, CALR, or MPL mutations. Presents with anemia, massive splenomegaly, constitutional symptoms (fever, night sweats, weight loss), and leukoerythroblastic blood smear. Diagnosis: bone marrow biopsy (fibrosis, megakaryocytic hyperplasia), JAK2 V617F mutation (50-60\%), CALR (20-25\%). Prognosis: dynamic IPSS. Treatment: JAK1/2 inhibitors (ruxolitinib, fedratinib), splenectomy for symptomatic splenomegaly, and allogeneic stem cell transplant as only curative option.

\medterm{Non-Hodgkin Lymphoma}-- Hematologic malignancy of B or T/NK cells. B-cell NHL (85\%) includes diffuse large B-cell lymphoma (DLBCL, most common aggressive), follicular lymphoma (most common indolent), mantle cell lymphoma, marginal zone lymphoma, and Burkitt lymphoma. T-cell NHL: peripheral T-cell lymphoma, anaplastic large cell lymphoma. Presents with painless lymphadenopathy, B symptoms (fever, night sweats, weight loss), and cytopenias. Diagnosis: excisional biopsy with flow cytometry, cytogenetics, FISH. Treatment: R-CHOP for DLBCL, rituximab-based therapy for follicular, and stem cell transplant for relapsed/refractory disease.

\medterm{Paroxysmal Nocturnal Hemoglobinuria}-- A rare, acquired clonal disorder of hematopoietic stem cells characterized by complement-mediated intravascular hemolysis, thrombosis (particularly in unusual sites such as hepatic, portal, and cerebral veins), and bone marrow failure. The underlying defect is a somatic mutation in the PIGA gene, which encodes a subunit of the GPI-anchor synthesis enzyme complex, resulting in deficient expression of GPI-anchored proteins (CD55 and CD59) on the surface of blood cells. Diagnosis: flow cytometry showing deficient CD55/CD59 on red blood cells. Treatment: eculizumab or ravulizumab (anti-C5 monoclonal antibodies) for hemolysis and thrombosis; anticoagulation for thrombotic events; supportive care with folic acid, transfusion, and iron chelation for transfusion burden.

\medterm{Pernicious Anemia}-- A type of vitamin B12 deficiency anemia caused by autoimmune destruction of gastric parietal cells, leading to lack of intrinsic factor and impaired B12 absorption in the terminal ileum. More common in older adults, those with autoimmune disorders, and pernicious anemia is associated with atrophic gastritis and increased gastric cancer risk. Presents with megaloblastic anemia (macrocytic RBCs, hypersegmented neutrophils), neuropathy (paresthesias, ataxia, cognitive impairment), and glossitis. Diagnosis: low B12, elevated methylmalonic acid and homocysteine, positive anti-intrinsic factor antibodies. Treatment: lifelong parenteral B12 (cyanocobalamin) replacement.

\medterm{Platelet Function Disorders}-- Inherited: Glanzmann thrombasthenia (GPIIb/IIIa deficiency, absent platelet aggregation, mucocutaneous bleeding, recombinant FVIIa, platelet transfusions), Bernard-Soulier syndrome (GPIb/IX/V deficiency, giant platelets, thrombocytopenia, prolonged bleeding time, platelet transfusions), Storage pool deficiency (delta/dense granules, mild bleeding). Acquired: uremia (platelet dysfunction, desmopressin, dialysis), myeloproliferative neoplasms (acquired von Willebrand syndrome in ET/PV, aspirin), myelodysplasia. Diagnosis: bleeding time, PFA-100, platelet aggregation studies, flow cytometry for glycoprotein expression. Treatment: desmopressin for uraemia and mild type 1 VWD; platelet transfusions for Glanzmann and Bernard-Soulier; recombinant FVIIa for Glanzmann with anti-GPIIb/IIIa antibodies.

\medterm{Platelets (Thrombocytes)}-- Small, anucleate, disc-shaped cell fragments (2--4 \(\mu\)m diameter) derived from megakaryocytes in the bone marrow, circulating at 150--400 \(\times\) 10\(^9\)/L. Platelets are essential for primary hemostasis: they adhere to exposed subendothelial collagen via von Willebrand factor (GPIb receptor), become activated (shape change, degranulation—release ADP, serotonin, thromboxane A$_2$), and aggregate via fibrinogen bridges (GPIIb/IIIa receptor) to form a platelet plug. They also provide a phospholipid surface for coagulation factor assembly (tenase and prothrombinase complexes). Platelet lifespan: 7--10 days; old platelets are removed by the spleen. Thrombocytopenia (<150 \(\times\) 10\(^9\)/L) increases bleeding risk; thrombocytosis (>450 \(\times\) 10\(^9\)/L) is reactive (infection, inflammation, iron deficiency) or clonal (essential thrombocythaemia). Platelet function disorders: Glanzmann thrombasthenia (GPIIb/IIIa deficiency), Bernard–Soulier syndrome (GPIb deficiency), storage pool disease. Drugs affecting platelets: aspirin (irreversible COX-1 inhibition), clopidogrel (P2Y\(_{12}\) receptor antagonist), NSAIDs, and GPIIb/IIIa inhibitors.

\medterm{Polycythemia Vera}-- A chronic myeloproliferative neoplasm characterized by increased red blood cell mass, often with increased white blood cell and platelet counts, driven by constitutive activation of the JAK-STAT signaling pathway. The JAK2 V617F mutation is present in approximately 95 percent of PV patients, with an additional 3 percent harboring JAK2 exon 12 mutations. Clinical presentation includes elevated hemoglobin and hematocrit (per WHO criteria: male hemoglobin greater than 16.5 g/dL or haematocrit greater than 49\% in men, greater than 16.0 g/dL or haematocrit greater than 48\% in women), thrombocytosis, leukocytosis, and splenomegaly in 70\% of patients. Treatment: phlebotomy (target Hct <45\%), low-dose aspirin for all patients, cytoreduction (hydroxyurea, interferon-alpha, ruxolitinib) for high-risk patients (age >60, prior thrombosis). Risk of thrombosis, hemorrhage, and transformation to myelofibrosis or AML.

\medterm{Primary Myelofibrosis}-- A chronic myeloproliferative neoplasm characterized by bone marrow fibrosis (reticulin and/or collagen fibrosis), extramedullary hematopoiesis (primarily splenic and hepatic), and a leukoerythroblastic peripheral blood picture. It results from clonal proliferation of hematopoietic stem cells carrying JAK2 V617F (50-60 percent), CALR (25-30 percent), or MPL (3-5 percent) mutations, leading to cytokine release that stimulates fibroblast proliferation and collagen deposition of extracellular matrix proteins. Clinical presentation includes progressive anemia, massive splenomegaly with early satiety and abdominal discomfort, constitutional symptoms (fever, night sweats, weight loss, fatigue), and extramedullary hematopoiesis. Peripheral blood shows leukoerythroblastosis with teardrop cells. Treatment: ruxolitinib (JAK1/2 inhibitor) for symptomatic splenomegaly and constitutional symptoms; allogeneic HSCT is the only potentially curative option; supportive care with transfusions, erythropoietin, and targeted therapy for iron overload.

\medterm{Protein C Deficiency}-- An inherited (autosomal dominant) thrombophilia causing increased VTE risk. Protein C is a vitamin K-dependent natural anticoagulant that inactivates factors Va and VIIIa and stimulates fibrinolysis. Heterozygous: 6-10 fold increased VTE risk. Homozygous: severe neonatal purpura fulminans (widespread skin necrosis and thrombosis) requiring protein C replacement. Diagnosis: functional and antigenic protein C levels (low). Caution: levels also low in warfarin therapy, liver disease, DIC, and vitamin K deficiency. Management: anticoagulation for VTE; warfarin initiation requires heparin bridging to avoid warfarin-induced skin necrosis.

\medterm{Protein S Deficiency}-- An inherited (autosomal dominant) thrombophilia. Protein S is a vitamin K-dependent natural anticoagulant that acts as a cofactor for activated protein C in inactivating factors Va and VIIIa. Heterozygous: 2-10 fold increased VTE risk. Homozygous: rare, purpura fulminans. Total, free, and functional protein S levels must be interpreted carefully as levels are affected by pregnancy, oral contraceptives, liver disease, DIC, and acute thrombosis. Diagnosis: confirmed low free protein S on repeat testing. Management: anticoagulation for VTE; lifelong anticoagulation for recurrent or severe thrombosis.

\medterm{Rhesus (Rh) Factor}-- A system of blood group antigens (the Rh system) on the surface of red blood cells, the most important of which is the D antigen (RhD). Individuals are Rh-positive (RhD+) if the D antigen is present (85\% of Caucasians, 92\% of Africans, 99\% of Asians) or Rh-negative (RhD-) if absent. The Rh factor is determined by the RHD gene on chromosome 1. Rh incompatibility occurs when an Rh-negative mother carries an Rh-positive fetus: fetal red cells entering the maternal circulation (fetomaternal hemorrhage—typically at delivery, but also during pregnancy, trauma, or abortion) stimulate maternal anti-D antibody production (IgG). In a subsequent Rh-positive pregnancy, maternal anti-D IgG crosses the placenta and causes hemolytic disease of the fetus and newborn (HDFN—erythroblastosis fetalis), leading to fetal anemia, hydrops fetalis, kernicterus, and neonatal jaundice. Prevention: prophylactic anti-D immunoglobulin (RhoGAM) is given to Rh-negative women at 28 weeks gestation and within 72 hours of delivery (or any sensitising event—amniocentesis, miscarriage, abortion, ectopic pregnancy), which coats fetal Rh-positive red cells and prevents maternal immune sensitisation. Other Rh antigens: C, c, E, e (Kell, Duffy, Kidd, and MNS are other clinically important blood group systems).

\medterm{Rh Factor}-- See Rhesus (Rh) Factor.

\medterm{Sickle Cell Trait}-- Heterozygous carrier state for the HbS mutation (HbAS), present in ~8\% of African Americans. Red cells contain both HbA and HbS (typically 30-40\% HbS). Usually asymptomatic with normal hemoglobin and RBC indices. Complications (rare but recognized): painless hematuria (renal papillary necrosis), urinary concentrating defect, splenic infarction at high altitude, exercise-related sudden death (during extreme exertion), and hyposthenuria. Diagnosis: hemoglobin electrophoresis showing HbA and HbS. Genetic counseling. No treatment needed; avoid extreme dehydration and high-altitude exposure.

\medterm{Splenomegaly}-- Enlargement of the spleen beyond its normal size (~12 cm, weight 150--200 g). Splenomegaly can be detected by percussion (Castell sign—dullness in the lowest intercostal space on inspiration) and palpation (normal spleen is not palpable; a palpable spleen indicates at least 2--3 times enlargement). Causes: congestive (portal hypertension—cirrhosis, hepatic vein thrombosis/Budd–Chiari syndrome, splenic vein thrombosis, heart failure), infectious (EBV, CMV, malaria, leishmaniasis, schistosomiasis, tuberculosis, endocarditis, brucellosis, HIV), hematological (hemolytic anemias—hereditary spherocytosis, thalassemia, sickle cell disease; myeloproliferative neoplasms—myelofibrosis, polycythemia vera, CML; leukemia—CLL, ALL, AML; lymphoma—Hodgkin, non-Hodgkin), inflammatory/autoimmune (sarcoidosis, SLE, rheumatoid arthritis—Felty syndrome, Still disease), infiltrative (Gaucher disease, Niemann–Pick, amyloidosis, glycogen storage diseases), and neoplastic (primary splenic tumors—rare). Massive splenomegaly (spleen >20 cm, >1000 g): most commonly myelofibrosis, CML, malaria, kala-azar, Gaucher disease, and schistosomiasis. Complications: hypersplenism (pancytopenia from increased sequestration/destruction of blood cells), splenic infarction (pain, left upper quadrant), and splenic rupture (trauma). Diagnosis: ultrasound, CT, MRI; further workup directed at underlying cause. Treatment: directed at the underlying cause; splenectomy for symptomatic massive splenomegaly, hypersplenism with severe cytopenias, splenic abscess, or splenic rupture.

\medterm{Stem Cell Transplantation}-- Hematopoietic stem cell transplantation (HSCT) replaces diseased or ablated bone marrow with healthy donor or autologous stem cells. Autologous HSCT: patient's own stem cells collected after mobilization (G-CSF with or without plerixafor), providing hematopoietic rescue after high-dose chemotherapy without graft-versus-host disease risk; indications include relapsed lymphoma and multiple myeloma. Allogeneic HSCT: donor stem cells (HLA-matched sibling, haploidentical family member, or matched unrelated donor). Conditioning: myeloablative (high-dose chemo + TBI) or reduced-intensity (older, comorbidities). GVHD prophylaxis: calcineurin inhibitor plus methotrexate or mycophenolate. Complications: GVHD (acute: skin, liver, GI; chronic: skin, mouth, eyes, liver, lungs), infection (CMV, EBV, fungal), VOD/SOS (veno-occlusive disease), graft failure, and relapse.

\medterm{T-Cell Lymphomas}-- Peripheral T-cell lymphomas (PTCL) are a heterogeneous group of aggressive lymphomas including PTCL-NOS (not otherwise specified, most common), angioimmunoblastic T-cell lymphoma (AITL, T-follicular helper origin, RHOA mutations, dysimmune features), anaplastic large cell lymphoma (ALCL, ALK+ vs ALK-), enteropathy-associated T-cell lymphoma (EATL, celiac association), hepatosplenic T-cell lymphoma (gamma-delta, isochromosome 7q), and cutaneous T-cell lymphomas (mycosis fungoides, Sezary syndrome (erythroderma, generalised lymphadenopathy, circulating Sezary cells). Diagnosis: excisional lymph node biopsy with histology, IHC, flow cytometry, cytogenetics, and FISH. Treatment: CHOP-like chemotherapy for aggressive subtypes; novel agents (brentuximab vedotin for ALCL; romidepsin/belinostat for relapsed PTCL); allogeneic HSCT for eligible patients.

\medterm{thalassemia}-- A group of inherited hemoglobin disorders caused by reduced or absent synthesis of one or more globin chains, leading to microcytic hypochromic anemia. Types: $\alpha$-thalassemia (deletion of $\alpha$-globin genes on chromosome 16; severity depends on number of deleted genes: 1 deletion—silent carrier, 2—$\alpha$-thalassemia trait/minor, 3—hemoglobin H disease (HbH—moderate hemolytic anemia), 4—Hb Bart hydrops fetalis (lethal in utero)), $\beta$-thalassemia (point mutations in $\beta$-globin gene on chromosome 11; $\beta$-thalassemia minor/trait—mild anemia, $\beta$-thalassemia intermedia—moderate anemia requiring intermittent transfusions, $\beta$-thalassemia major (Cooley anemia)—severe anemia presenting at 6--24 months, requiring lifelong transfusions, leads to iron overload, growth retardation, hepatosplenomegaly, and skeletal changes (bossing, chipmunk facies, osteopenia)). Epidemiology: high prevalence in the Mediterranean, Middle East, South Asia, and Southeast Asia (malaria protection). Diagnosis: full blood count (microcytosis, hypochromia), hemoglobin electrophoresis (HbA$_2$ elevated in $\beta$-thalassemia trait, HbF elevated in $\beta$-thalassemia major, HbH in $\alpha$-thalassemia intermedia), and genetic testing. Management: mild forms require no treatment, or folate supplementation. Transfusion-dependent patients: regular packed red cell transfusions (maintain Hb 9--10 g/dL), iron chelation therapy (deferasirox, desferrioxamine, or deferiprone) to prevent iron overload—the main cause of morbidity (cardiac, hepatic, endocrine). Allogeneic hematopoietic stem cell transplantation is the only curative option (best outcomes in children <16 years without significant iron overload). Luspatercept (erythroid maturation agent) reduces transfusion burden in $\beta$-thalassemia. Gene therapy (betibeglogene autotemcel—Zynteglo) is approved for transfusion-dependent $\beta$-thalassemia.

\medterm{Thrombectomy}-- A surgical or interventional radiology procedure to remove a blood clot (thrombus) from a blood vessel, restoring blood flow. Types: (1) Endovascular thrombectomy (mechanical thrombectomy)—for acute ischemic stroke due to large vessel occlusion (LVO—internal carotid artery, middle cerebral artery M1 segment). Performed via femoral artery access, a stent retriever or aspiration catheter is deployed to remove the clot. Window: up to 6--24 hours from symptom onset (selected by perfusion imaging). Number needed to treat (NNT) of 3 for functional independence. (2) Surgical thrombectomy—open arteriotomy and clot removal, rarely used now (reserved for failed endovascular therapy, massive PE with hemodynamic instability, or peripheral arterial occlusion). (3) Venous thrombectomy—for extensive iliofemoral DVT (phlegmasia cerulea dolens) or superior vena cava syndrome. (4) Vacuum-assisted thrombectomy (AngioJet, Penumbra)—for peripheral arterial or venous thrombosis. Complications: vessel perforation, dissection, distal embolisation, hemorrhage (especially intracranial hemorrhage post-thrombectomy for stroke).

\medterm{Thrombocytes}-- See Platelets (Thrombocytes).

\medterm{Thrombophilia}-- A tendency to develop venous thromboembolism (VTE) due to inherited or acquired hypercoagulable states. Inherited: factor V Leiden, prothrombin G20210A mutation, antithrombin deficiency, protein C deficiency, protein S deficiency. Acquired: antiphospholipid syndrome (lupus anticoagulant, anticardiolipin antibodies, beta-2 glycoprotein I antibodies), cancer, pregnancy, oral contraceptives, and obesity. Presents with DVT and PE, often recurrent or at unusual sites. Management: anticoagulation (DOACs, warfarin, heparin), with duration depending on provoking factors and risk of recurrence.

\medterm{Thrombophilia Testing}-- Thrombophilia testing evaluates for inherited and acquired hypercoagulable states in patients with venous thromboembolism, recurrent pregnancy loss, or strong family history of VTE. Inherited thrombophilias include Factor V Leiden (most common, 5 percent of Caucasians, activated protein C resistance, detected by clot-based or genetic assay), prothrombin G20210A mutation (2 percent of Caucasians, elevated prothrombin levels), antithrombin deficiency (1 percent, most thrombogenic, measured by activity (functional assay), protein C deficiency (activity assay), protein S deficiency (free protein S antigen, activity). Acquired thrombophilias: antiphospholipid syndrome (lupus anticoagulant, anticardiolipin, anti-beta-2 glycoprotein I), malignancy-associated thrombosis. Testing should be performed at least 3 months after a thrombotic event and after discontinuation of anticoagulation to avoid false-positive or false-negative results.

\medterm{Thrombotic Thrombocytopenic Purpura}-- A life-threatening thrombotic microangiopathy caused by severe deficiency of ADAMTS13 (a disintegrin and metalloproteinase with thrombospondin type 1 motif member 13), a von Willebrand factor-cleaving protease. ADAMTS13 deficiency (activity less than 10 percent) leads to accumulation of ultra-large von Willebrand factor multimers that cause spontaneous platelet adhesion and aggregation in the microcirculation, forming platelet-rich thrombi that consume platelets and coagulation factors, causing thrombocytopenia and microangiopathic hemolytic anemia (schistocytes on peripheral smear). Clinical: fever, thrombocytopenia, microangiopathic hemolytic anemia, renal impairment, neurologic symptoms (headache, confusion, seizure, stroke), and fever. Diagnosis: ADAMTS13 activity <10\% and detectable anti-ADAMTS13 antibodies. Treatment: daily plasma exchange (PLEX) until platelet count normalizes and hemolysis resolves; corticosteroids for immune suppression; rituximab for refractory/relapsed TTP; caplacizumab (anti-vWF nanobody) for acute TTP.

\medterm{Transfusion Medicine}-- Blood components: PRBCs (leukoreduced, irradiated for at-risk), FFP (coagulation factors), platelets (apheresis pooled, leukoreduced), cryoprecipitate (fibrinogen, vWF, factor VIII, XIII). ABO/Rh compatibility. Transfusion thresholds: restrictive (Hb 7-8 g/dL) for most stable patients (TRICC, FOCUS, TRISS trials); liberal (Hb 9-10) for acute coronary syndrome, severe TBI, acute bleeding. Massive transfusion protocol (MTP): 1:1:1 (PRBC:FFP:platelets) or 1:1:2; TXA within 3h (CRASH-2). Adverse reactions: acute hemolytic (ABO incompatibility, fever, hypotension, DIC, renal failure), febrile non-hemolytic (most common, cytokines in stored blood), allergic (urticaria, anaphylaxis in IgA deficiency), transfusion-related acute lung injury (TRALI, bilateral pulmonary infiltrates, hypoxia), transfusion-associated circulatory overload (TACO, pulmonary edema from volume overload), and transfusion-transmitted infections (HIV, HBV, HCV, West Nile virus, bacterial contamination).

\medterm{Vitamin B12 Deficiency}-- Vitamin B12 (cobalamin) deficiency causes megaloblastic anemia and neurologic dysfunction through impaired DNA synthesis and abnormal fatty acid metabolism in myelin. Dietary sources include animal products (meat, fish, dairy), with daily requirement of 2-3 mcg. Absorption requires adequate gastric acid and pepsin to release B12 from food proteins, binding to intrinsic factor (IF) produced by gastric parietal cells, and absorption in the terminal ileum via the cubilin-amnionless receptor complex. Causes: pernicious anemia (autoimmune destruction of parietal cells, most common in adults), dietary deficiency (vegans, malnutrition, alcoholism), malabsorption (gastric bypass, Crohn disease, celiac disease, pancreatic insufficiency), terminal ileal disease or resection, and drugs (metformin, PPI). Lab: macrocytic anemia (MCV >100), hypersegmented neutrophils, low B12, elevated methylmalonic acid and homocysteine. Treatment: intramuscular cyanocobalamin 1000 mcg daily for one week, then weekly for one month, then monthly for life; high-dose oral B12 (1000-2000 mcg/day) may be adequate for dietary deficiency.

\medterm{von Willebrand Disease}-- The most common inherited bleeding disorder, caused by quantitative or qualitative deficiency of von Willebrand factor (VWF), a multimeric glycoprotein that mediates platelet adhesion to subendothelium and carries factor VIII. Type 1 (partial quantitative deficiency, 70-80\%): mild mucocutaneous bleeding, easy bruising, menorrhagia. Type 2 (qualitative defects): 2A (loss of high-molecular-weight multimers, decreased platelet adhesion), 2B (increased affinity for platelet GPIb (thrombocytopenia), 2M (normal multimers, decreased platelet adhesion), 2N (decreased factor VIII binding, mimics hemophilia A). Type 3 (severe quantitative deficiency). Treatment: desmopressin (DDAVP) for mild type 1; von Willebrand factor-containing concentrates (Humate-P) for more severe bleeding or surgical prophylaxis; tranexamic acid for mucocutaneous bleeding; avoid aspirin and NSAIDs.

\medterm{Waldenstrom Macroglobulinemia}-- A rare B-cell lymphoplasmacytic lymphoma characterized by bone marrow infiltration by lymphoplasmacytic cells and production of monoclonal IgM paraprotein. It accounts for less than 2 percent of hematologic malignancies and typically affects older adults (median age 63-65 years). Clinical manifestations result from both marrow infiltration (anemia, cytopenias) and IgM-related effects: hyperviscosity syndrome (visual disturbances, neurologic symptoms, bleeding, and cryoglobulinaemia. Diagnosis: serum protein electrophoresis (M-spike), immunofixation (IgM), bone marrow biopsy (lymphoplasmacytic infiltrate >10\%). Treatment: symptomatic or progressive disease: rituximab alone or in combination with bendamustine, cyclophosphamide/dexamethasone, or bortezomib; plasmapheresis for hyperviscosity syndrome; ibrutinib for MYD88-mutated disease; venetoclax for relapsed/refractory.

\medterm{Warfarin Management}-- Warfarin (vitamin K antagonist) inhibits gamma-carboxylation of factors II, VII, IX, X, proteins C and S. Target INR: 2-3 (most indications), 2.5-3.5 (mechanical mitral valve, recurrent VTE on therapeutic INR). Initiation: 5 mg daily x2-3 days, then adjust per INR; genotype-guided (CYP2C9, VKORC1) algorithms. Interactions: numerous (antibiotics, amiodarone, antifungals, NSAIDs, herbal). Monitoring: INR frequency per stability. Bridging: therapeutic LMWH when INR subtherapeutic for high thrombotic risk. Management: VTE treatment with DOACs or LMWH bridging to warfarin. Long-term anticoagulation based on recurrence risk and bleeding risk. Regular INR monitoring, dietary consistency (vitamin K intake), and medication review to avoid interactions.

\medterm{White Cell}-- Also known as leukocytes, the cells of the immune system that defend the body against infection and foreign substances. Produced in the bone marrow from hematopoietic stem cells. Types: neutrophils (50-70\%, first responders to bacterial infection, phagocytosis), lymphocytes (20-40\%, adaptive immunity — T-cells, B-cells, NK cells), monocytes (2-8\%, differentiate into macrophages and dendritic cells), eosinophils (1-4\%, allergic reactions, parasitic infections), and basophils (<1\%, allergic responses, release histamine). Normal WBC count: 4,000-11,000/\(\mu\)L. Leukocytosis (elevated WBC): infection, inflammation, stress, leukemia. Leukopenia (low WBC): bone marrow suppression, infection, autoimmune disorders.

\medterm{White Cell Count}-- A laboratory measurement of the number of white blood cells (leukocytes) in a given volume of blood, part of the complete blood count (CBC). Normal range: 4,000-11,000 cells/\(\mu\)L (4-11 × 10\(^9\)/L). Differential white cell count measures the proportion of each leukocyte type: neutrophils, lymphocytes, monocytes, eosinophils, basophils. An elevated WBC (leukocytosis) suggests infection, inflammation, tissue damage, or hematologic malignancy. A decreased WBC (leukopenia) suggests bone marrow suppression from chemotherapy, infection (HIV, sepsis), autoimmune destruction, or drug toxicity. Marked leukocytosis (>100,000/\(\mu\)L) is a defining feature of chronic leukemias (CLL, CML). Left shift: increased band neutrophils, indicating acute bacterial infection.

